

\section{\texorpdfstring{$\lieSO(1, k)$}{so(1,k)}} \label{section:TypeA1}

In this section, we deal with the $A_1$ case defined in Definition \ref{def:typeA}.
We use the theory of Jordan algebras and $3$-graded Lie algebras to study the structure of $\lie{g}_1$.

\subsection{Jordan algebras and graded Lie algebras} \label{subsect:JordanAlgebra}

We shall study the case of type $A_1$ defined in Definition \ref{def:typeA}.
To do this, we first recall basic properties of Jordan algebras and $3$-graded Lie algebras.
We refer the reader to \cite{FaKo94} and \cite{HiKoMo14} for details of the relation between Jordan algebras and $3$-graded Lie algebras.
See also the references therein.

Let $\lie{g} = \lie{g}_1 \oplus \lie{g}_0 \oplus \lie{g}_{-1}$ be a graded simple Lie algebra with Cartan involution $\theta$ such that $\theta(\lie{g}_i) = \lie{g}_{-i}$, and $h \in \lie{g}_0$ the characteristic element of the grading.
We equip $\lie{g}_1$ with a Jordan triple product by $\jprod{x, y, z} \coloneq \frac{1}{2}[[x, -\theta(y)], z]$.
The product satisfies the identities:
\begin{align*}
	\jprod{x, y, z} &= \jprod{z, y, x}, \\
	\jprod{x, y, \jprod{u, v, w}} &= \jprod{\jprod{x, y, u}, v, w} - \jprod{u, \jprod{y, x, v}, w} + \jprod{u, v, \jprod{x, y, w}}
\end{align*}
for any $x, y, z, u, v, w \in \lie{g}_1$.

We say that $x \in \lie{g}_1$ is a \define{tripotent} if $\jprod{x, x, x} = x$.
In this paper, by a tripotent we always mean a nonzero tripotent.
The following proposition is clear from the definition of the Jordan triple product.

\begin{proposition} \label{prop:TripotentAndSL2Triple}
	Let $x \in \lie{g}_1$.
	Then $\set{[x, -\theta(x)], x, -\theta(x)}$ forms an $\lieSL_2$-triple if and only if $x$ is a tripotent.
\end{proposition}

Assume that there exists an element $u \in \lie{g}_1$ such that $[u, -\theta(u)] = 2h$.
Such an element is called a \define{unitary} element.
Then we have $\jprod{u, u, x} = x$ for any $x \in \lie{g}_1$ and, in particular, $u$ is a tripotent.
This implies that $\lie{g}_1$ has a Jordan algebra structure with the product $x \circ y = \jprod{x, u, y}$ and unit $u$, that is, the product $\circ$ satisfies the identities:
\begin{align*}
	x\circ y = y\circ x, \quad x\circ (x^2 \circ y) = x^2 \circ (x\circ y)
\end{align*}
for any $x, y \in \lie{g}_1$, where $x^2 = x\circ x$.
Note that the Jordan algebra structure depends on the choice of $u$.
When we need to specify the unit, we denote by $(\lie{g}_1, u)$ the Jordan algebra $\lie{g}_1$ with the unit $u$.
For a while, we fix the Jordan algebra structure on $\lie{g}_1$.

We equip $\lie{g}_1$ with a trace and a symmetric bilinear form by
\begin{align*}
	\tr(x) \coloneq -(x, \theta(u)), \tau(x, y) \coloneq \tr(x \circ y).
\end{align*}
Note that $\tau(\cdot, \cdot)$ and $\tr$ are scalar multiples of the standard bilinear form and trace, respectively, on the Jordan algebra $\lie{g}_1$.
The scalar multiples are not important in this paper, so we do not specify them.
Define $\vartheta(x) \coloneq \jprod{u, x, u}$ for $x \in \lie{g}_1$.
Then $\vartheta$ is an involution of the Jordan algebra $\lie{g}_1$ such that
\begin{align*}
	\tau(x, \vartheta(y)) = \tau(\vartheta(x), y) = -(x, \theta(y)) \quad (\forall x, y \in \lie{g}_1).
\end{align*}
Hence $\tau(\cdot, \vartheta(\cdot))$ is an inner product on $\lie{g}_1$.

\begin{proposition} \label{prop:CartanInvOrthogonal}
	One has $\lie{g}_1^{\vartheta} \perp \lie{g}_1^{-\vartheta}$ with respect to the inner product $\tau(\cdot, \vartheta(\cdot)) = -(\cdot, \theta(\cdot))$.
\end{proposition}

We call $\vartheta$ the \define{Cartan involution} of the Jordan algebra $\lie{g}_1$.
The triple product $\jprod{\cdot, \cdot, \cdot}$ is reconstructed from the Jordan product $\circ$ and the Cartan involution $\vartheta$ by
\begin{align}
	\jprod{x, y, z} = (x \circ \vartheta(y)) \circ z + (z \circ \vartheta(y)) \circ x - (x \circ z) \circ \vartheta(y). \label{eqn:TripleProductFromJordanProduct}
\end{align}


To study orbit decompositions in $\lie{g}_1$, we use Jordan frames.
An element $x \in \lie{g}_1$ is called an \define{idempotent} if $x \circ x = x$ and $x \neq 0$.
For two idempotents $x, y \in \lie{g}_1$, we say that $x$ and $y$ are \define{orthogonal} if $x \circ y = 0$.
For an idempotent $x$, if there exist no orthogonal idempotents $x_1, x_2$ with $x = x_1 + x_2$, then $x$ is called \define{primitive}.
A family of mutually orthogonal primitive idempotents $\set{u_1, \ldots, u_r}$ in $\lie{g}_1^{\vartheta}$ is called a \define{Jordan frame} if $u = \sum_{i=1}^r u_i$ and $r$ is as large as possible.
We rephrase the notions in terms of the Lie algebra $\lie{g}$.

\begin{proposition} \label{prop:IdempotentAndOrthogonal}
	Let $x, y \in \lie{g}_1^\vartheta$ be idempotents.
	\begin{enumparen}
		\item $\set{[x, -\theta(x)], x, -\theta(x)}$ forms an $\lieSL_2$-triple.
		\item $x$ and $y$ are orthogonal if and only if $[x, -\theta(y)] = [y, -\theta(x)] = 0$.
	\end{enumparen}
\end{proposition}

\begin{proof}
	(1) This follows from Proposition \ref{prop:TripotentAndSL2Triple} and \eqref{eqn:TripleProductFromJordanProduct}.

	(2) Assume that $x$ and $y$ are orthogonal.
	For $a \in \lie{g}_1$, let $L(a)$ denote the linear map $a\circ (\cdot)$.
	Then we have the Jordan identity $[L(a^2), L(a)] = 0$.
	For any $a, b\in \RR$, expanding the left-hand side of the identity $[L((ax+by)^2), L(ax + by)] = 0$,
	we have
	\begin{align*}
		a^2b[L(x), L(y)] + ab^2[L(y), L(x)] = 0.
	\end{align*}
	This implies that $[L(x), L(y)] = 0$.
	By \eqref{eqn:TripleProductFromJordanProduct}, we have
	\begin{align*}
		\jprod{x, y, z} = (z \circ y) \circ x - (x \circ z) \circ y = L(x)L(y)z - L(y)L(x)z = 0
	\end{align*}
	for any $z \in \lie{g}_1$.
	This shows that $[x, -\theta(y)] =  [y, -\theta(x)] = 0$.

	To show the converse, assume that $[x, -\theta(y)] = [y, -\theta(x)] = 0$.
	By \eqref{eqn:TripleProductFromJordanProduct}, we have
	\begin{align*}
		x\circ y = (x \circ y) \circ u = \frac{1}{2}\left(\jprod{x, y, u} + \jprod{y, x, u}\right) = 0.
	\end{align*}
	We have shown the assertion.
\end{proof}

The following fact is a consequence of the Peirce decomposition \cite[Theorem IV.2.1]{FaKo94} of the Jordan algebra $\lie{g}_1$ with respect to a Jordan frame.
Let $\sigma$ be the involution of $\lie{g}$ such that $\lie{g}^{\sigma} = \lie{g}_0$.
Let $K_0$ be the analytic subgroup of $G$ with the Lie algebra $\lie{g}^\theta_0 = \lie{k}^\sigma$.


\begin{fact} \label{fact:JordanFrame}
	Let $\set{u_1, \ldots, u_r}$ be a Jordan frame in $\lie{g}_1^{\vartheta}$.
	Set $h_i \coloneq [u_i, -\theta(u_i)]$ for $i = 1, \ldots, r$.
	\begin{enumparen}
		\item $[u_i, -\theta(u_j)] = 0$ for any $i \neq j$.
		\item $[h_i, h_j] = 0$ for any $i, j$.
		\item $[h_i, u_j] = 2\delta_{ij} u_j$ for any $i, j$.
		\item $\sum_{i=1}^r h_i = 2h$.
		\item $\spn{\RR}{u_1 - \theta(u_1), \ldots, u_{r} - \theta(u_{r})}$ is a maximal abelian subspace in $\lie{g}^{-\theta, -\sigma} = (\lie{g}_1 + \lie{g}_{-1})\cap \lie{g}^{-\theta}$.
		\item There exists a maximal abelian subspace $\lie{a}$ of $\lie{g}_0^{-\theta}$ such that $u_i$ are root vectors and the roots are strongly orthogonal.
	\end{enumparen}
\end{fact}

\begin{proposition} \label{prop:K0Orbit}
	Let $\set{u_1, \ldots, u_r}$ be a Jordan frame in $\lie{g}_1^{\vartheta}$, and set $\lie{c}\coloneq \spn{\RR}{u_1, \ldots, u_r}$.
	One has $\Ad(K_0)\lie{c} = \lie{g}_1$.
\end{proposition}

\begin{proof}
	Define $\varphi\colon \lie{g}_1 \rightarrow \lie{g}^{-\theta, -\sigma}$ by $\varphi(x) = x -\theta(x)$.
	Then $\varphi$ is a $K_0$-equivariant linear isomorphism.
	By Fact \ref{fact:JordanFrame} (5), $\varphi(\lie{c})$ is a maximal abelian subspace in $\lie{g}^{-\theta, -\sigma}$.
	This implies that $\Ad(K_0)\varphi(\lie{c}) = \lie{g}^{-\theta, -\sigma} = \varphi(\lie{g}_1)$.
	Since $\varphi$ is bijective, we obtain $\Ad(K_0)\lie{c} = \lie{g}_1$.
\end{proof}

\begin{proposition} \label{prop:TripotentAndJordanFrame}
	Let $\set{u_1, \ldots, u_r}$ be a Jordan frame in $\lie{g}_1^{\vartheta}$.
	Let $v \in \lie{g}_1$ be a tripotent.
	Then there exist an element $k \in K_0$ and $\varepsilon_1, \ldots, \varepsilon_r \in \{-1, 0, 1\}$ such that $\Ad(k)v = \sum_{i=1}^r \varepsilon_i u_i$.
	Moreover, if $[v, -\theta(v)] = 2h$, then $\varepsilon_i \in \{-1, 1\}$ for any $i$.
\end{proposition}

\begin{proof}
	By Proposition \ref{prop:K0Orbit}, there exist $k \in K_0$ and $\varepsilon_1, \ldots, \varepsilon_r \in \RR$ such that
	\begin{align*}
		\Ad(k)v = \sum_{i=1}^r \varepsilon_i u_i.
	\end{align*}
	Since $\jprod{v, v, v} = v$, we have $\sum_{i=1}^r \varepsilon_i^3 u_i = \sum_{i=1}^r \varepsilon_i u_i$.
	This implies that $\varepsilon_i \in \{-1, 0, 1\}$ for any $i$.
	The last assertion follows from a similar calculation.
\end{proof}

Let $G_0$ be the analytic subgroup of $G$ with the Lie algebra $\lie{g}_0$.
We shall review the orbit decomposition of $\lie{g}_1$ under the action of $G_0$.
We refer the reader to \cite{Ka98} for details.

Assume that $\lie{g}$ is simple.
By Table \ref{table:JordanAlgebra} (see \cite[Table II]{Ka98} and \cite[Table 4]{HiKoMo14}), there exists $u \in \lie{g}_1$ such that $\lie{g}_1^\vartheta$ is a simple Jordan algebra.
Note that there may be several possible choices of $u$ for the Jordan triple system $\RR^{p,q}$ up to the $K_0$-action, i.e., $\RR^{p,q} \simeq \RR^{q,p}$ as Jordan triple systems.
In the table, we fix the unit vector $e_1$ as $u$.
We treat the case of $\lie{g} \simeq \lieSO(p, q)$ in Subsection \ref{subsect:JordanAlgebraSO(p,q)}.

\begin{table}[htbp]
	\centering
	\small
	\renewcommand{\arraystretch}{1.15}
	\begin{tabular}{c|c|c}
		$\lie{g}$ & $\Delta(\lie{g}^{\theta\sigma}, \lie{a}')$ & $\lie{g}_1$ \\
		\hline
		$\lieSL(2, \RR)$ & $A_0$ & $\RR$ \\
		$\lieSp(r, \RR)$ ($r \geq 3$) & $A_{r-1}$ & $\operatorname{Sym}(r, \RR)$ \\
		$\lieSU(r, r)$ ($r \geq 3$) & $A_{r-1}$ & $\operatorname{Herm}(r, \CC)$ \\
		$\lieSO^*(4r)$ ($r \geq 3$) & $A_{r-1}$ & $\operatorname{Herm}(r, \HH)$ \\
		$\lieSO(2, q+1)$ ($q \geq 2$) & $A_1$ & $\RR^{1,q}$ \\
		$\lieE_{7(-25)}$ & $A_2$ & $\operatorname{Herm}(3, \mathbb{O})$ \\ \hline
		$\lieSL(2r, \RR)$ ($r \geq 3$) & $D_r$ & $M(r, \RR)$ \\
		$\lieSO(2r, 2r)$ ($r \geq 3$) & $D_r$ & $\operatorname{Skew}(2r, \RR)$ \\
		$\lieSO(p+1, q+1)$ ($p\geq 2,q \geq 2$) & $D_2$ & $\RR^{p,q}$ \\
		$\lieE_{7(7)}$ & $D_3$ & $\operatorname{Herm}(3, \mathbb{O}_s)$ \\ \hline
		$\lieSp(r, \CC)$ ($r \geq 3$) & $C_r$ & $\operatorname{Sym}(r, \CC)$ \\
		$\lieSL(2r, \CC)$ ($r \geq 3$) & $C_r$ & $M(r, \CC)$ \\
		$\lieSO(4r, \CC)$ ($r \geq 3$) & $C_r$ & $\operatorname{Skew}(2r, \CC)$ \\
		$\lieSO(n+2, \CC)$ ($n \geq 3$) & $C_2$ & $\CC^n$ \\
		$\lieE_{7(\CC)}$ & $C_3$ & $\operatorname{Herm}(3, \mathbb{O})_{\CC}$ \\ \hline
		$\lieSp(r, r)$ ($r \geq 2$) & $C_r$ & $\operatorname{Sym}(2r, \CC) \cap M(r, \HH)$ \\
		$\lieSL(2r, \HH)$ ($r \geq 2$) & $C_r$ & $M(r, \HH)$ \\
		$\lieSO(p+1, 1)$ ($p \geq 2$) & $C_1$ & $\RR^{p,0}$
	\end{tabular}
	\caption{Classification of $3$-graded simple Lie algebras with unitary element $u$ and corresponding Jordan algebras}
	\label{table:JordanAlgebra}
\end{table}

Fix a Jordan frame $\set{u_1, \ldots, u_r}$ in $\lie{g}_1^\vartheta$.
Set
\begin{align}
	o_{p, q} &\coloneq \sum_{i=1}^p u_i - \sum_{j=p+1}^{p+q} u_j \quad (p, q \geq 0, p+q \leq r), \nonumber \\
	\lie{c} &\coloneq \spn{\RR}{u_1, \ldots, u_r}, \nonumber \\
	\lie{a'} &\coloneq \spn{\RR}{u_1 - \theta(u_1), \ldots, u_{r} - \theta(u_{r})}. \label{eqn:DefineOCA}
\end{align}
By Fact \ref{fact:JordanFrame} (5), $\lie{a'}$ is a maximal abelian subspace in $\lie{g}^{-\theta, -\sigma}$.

\begin{fact}[{\cite[Table II]{Ka98}}] \label{fact:RootSystemJordan}
	The root system $\Delta(\lie{g}^{\theta\sigma}, \lie{a'})$ is either of type $A_{r-1}$, $C_r$ or $D_r$.
	If the root system is of type $A_{r-1}$, then $\lie{g}^{\theta\sigma}$ has one-dimensional center $\RR(u - \theta(u))$.
\end{fact}

\begin{remark}
	When $\lie{g} \simeq \lieSL(2,\RR)$, the root system has no roots.
	We regard the root system as of type $A_0$ in this case.
	When $\lie{g} \simeq \lieSO(1, k)$ with $k \geq 3$, the root system is of type $A_1$.
	However, we regard the root system as of type $C_1$ for the results below.
\end{remark}

\begin{proposition} \label{prop:EuclideanJordanRoot}
	$\vartheta = \id$ if and only if $\Delta(\lie{g}^{\theta\sigma}, \lie{a'})$ is of type $A_{r-1}$.
\end{proposition}

\begin{remark}
	If the equivalent conditions in Proposition \ref{prop:EuclideanJordanRoot} hold, then $\lie{g}_1$ is called \define{Euclidean}.
\end{remark}

\begin{proof}
	Assume that $\Delta(\lie{g}^{\theta\sigma}, \lie{a'})$ is of type $A_{r-1}$.
	Then $u$ is a $K_0$-fixed vector and hence $\vartheta$ is $K_0$-equivariant.
	Let $x \in \lie{g}_1$.
	We shall show that $\vartheta(x) = \jprod{u, x, u} = x$.
	By Proposition \ref{prop:K0Orbit}, we may assume that $x \in \lie{c}$.
	Since $\jprod{u, u_i, u} = u_i$ for any $i$, we obtain $\jprod{u, x, u} = x$.

	Assume that $\vartheta = \id$.
	Let $x \in \lie{g}_1$.
	Then we have
	\begin{align*}
		\ad(u - \theta(u))^2 (x - \theta(x)) &= [u - \theta(u), [u,-\theta(x)] + [-\theta(u), x]] \\
		&= 2(-\vartheta(x) + x - \theta(x) + \theta(\vartheta(x))) = 0.
	\end{align*}
	This implies that $\ad(u - \theta(u))^2$ acts on $\lie{g}^{-\theta, -\sigma}$ trivially.
	Since $u - \theta(u)$ is a semisimple element in $\lie{g}^{-\theta, -\sigma}$, $u-\theta(u)$ belongs to the center of $\lie{g}^{\theta\sigma}$.
	By Fact \ref{fact:RootSystemJordan}, $\Delta(\lie{g}^{\theta\sigma}, \lie{a'})$ is of type $A_{r-1}$.
\end{proof}

\begin{fact}[{\cite{Ka98} (see also \cite[Part II, Theorem II.2.5]{FaKa00_jordan})}] \label{fact:JordanOrbitDecomposition}
	One has
	\begin{align*}
		\lie{g}_1 &= \coprod_{p, q \geq 0, p+q \leq r} \Ad(G_0)o_{p, q} && (\Delta(\lie{g}^{\theta\sigma}, \lie{a'}): A_{r-1}), \\
		\lie{g}_1 &= \coprod_{p=0}^r \Ad(G_0)o_{p, 0} && (\Delta(\lie{g}^{\theta\sigma}, \lie{a'}): C_r), \\
		\lie{g}_1 &= \coprod_{p=0}^r \Ad(G_0)o_{p, 0} \sqcup \Ad(G_0)o_{r-1, 1} && (\Delta(\lie{g}^{\theta\sigma}, \lie{a'}): D_r).
	\end{align*}
\end{fact}

Since the linear map $\lie{g}_1 \ni X\rightarrow X - \theta(X) \in \lie{g}^{-\theta, -\sigma}$ is $K_0$-equivariant, the Weyl group $W'$ of the restricted root system $\Delta(\lie{g}^{\theta\sigma}, \lie{a'})$ acts on $\lie{c}$.
The representation is isomorphic to the natural representation and the isomorphism is given by
\begin{align*}
	\lie{c} \ni \sum_{i=1}^r a_i u_i \mapsto (a_1, a_2, \ldots, a_r) \in \RR^r.
\end{align*}
For example, if $\Delta(\lie{g}^{\theta\sigma}, \lie{a'})$ is of type $A_{r-1}$, then $W'$ acts on the set $\set{u_1, \ldots, u_r}$ by permutations.
See the discussion following \cite[Lemma 3.2]{Ka98}, and also the discussion preceding \cite[Part II, Lemma II.2.2]{FaKa00_jordan}.

\begin{fact} \label{fact:WeylGroupJordan}
	The image of $W'$ in $\LieGL(\lie{c})$ coincides with the Weyl group of the root system:
	\begin{align*}
		&\set{\pm (u_i - u_j) : 1\leq i < j \leq r} && (\Delta(\lie{g}^{\theta\sigma}, \lie{a'}): A_{r-1}), \\
		&\set{\pm (u_i \pm u_j) : 1\leq i < j \leq r} \sqcup \set{\pm 2u_i: 1\leq i \leq r} && (\Delta(\lie{g}^{\theta\sigma}, \lie{a'}): C_r), \\
		&\set{\pm (u_i \pm u_j) : 1\leq i < j \leq r} && (\Delta(\lie{g}^{\theta\sigma}, \lie{a'}): D_{r}).
	\end{align*}
\end{fact}

We prove several lemmas that will be used in a later section to establish the non-minimality of NDD tuples.

\begin{lemma} \label{lem:RankOneTripotent}
	Let $X \in \Ad(G_0)o_{1, 0} \cup \Ad(G_0)o_{0, 1}$.
	Then there exists a constant $c > 0$ such that $cX$ is a tripotent, and hence $\set{[cX, -\theta(cX)], cX, -\theta(cX)}$ forms an $\lieSL_2$-triple.
\end{lemma}

\begin{proof}
	By Proposition \ref{prop:K0Orbit} and Fact \ref{fact:JordanOrbitDecomposition}, there exists $k\in K_0$ such that $\Ad(k)X \in \RR o_{1, 0} \cup \RR o_{0, 1}$.
	Since the property of being a tripotent is preserved under the action of $K_0$, the assertion follows.
\end{proof}


\begin{lemma} \label{lem:JordanAnyInnerProduct}
	Let $u' = o_{p,q}$ with $p+q = r$, and $X$ be a non-zero element in $\lie{g}_1$.
	If $\Delta(\lie{g}^{\theta\sigma}, \lie{a'})$ is of type $A_{r-1}$, assume that $p, q\neq 0$ in addition.
	Then, for each non-zero $a \in \RR$, there exists an element $x \in \overline{\Ad(G_0)X} \cap \lie{c}$ such that $(x, -\theta(u')) = a$.
\end{lemma}

\begin{proof}
	Assume that $\Delta(\lie{g}^{\theta\sigma}, \lie{a'})$ is not of type $A_{r-1}$.
	By Fact \ref{fact:JordanOrbitDecomposition}, $\RR o_{1, 0}$ is contained in $\overline{\Ad(G_0)X} \cap \lie{c}$.
	Hence some element in $\RR o_{1, 0}$ satisfies the required condition on $x$.

	Assume that $\Delta(\lie{g}^{\theta\sigma}, \lie{a'})$ is of type $A_{r-1}$.
	By Facts \ref{fact:JordanOrbitDecomposition} and \ref{fact:WeylGroupJordan}, $\varepsilon u_1$ and $\varepsilon u_r$ belong to $\overline{\Ad(G_0)X} \cap \lie{c}$ for some $\varepsilon \in \set{-1, 1}$.
	By assumption, we have $(u_1, -\theta(u')) > 0$ and $(u_r, -\theta(u')) < 0$.
	Hence some element in $\RR_{>0} \varepsilon u_1 \cup \RR_{>0} \varepsilon u_r$ satisfies the required condition on $x$.
\end{proof}

The following proposition is not necessary in this section.
It is used in the proof of Lemma \ref{lem:SUTopGrading}.

\begin{proposition} \label{prop:CodimensionOneJordan}
	Assume that $\dim(\lie{g}_1) \geq 2$.
	Let $W$ be a subspace of $\lie{g}_1$ such that $\Ad(G_0)W \cap \bigcup_{p+q=r} \Ad(G_0)o_{p,q} = \emptyset$.
	Then one has $\dim(W) < \dim(\lie{g}_1) - 1$.
\end{proposition}

\begin{proof}
	Assume that $\dim(W) = \dim(\lie{g}_1) - 1$.
	Set $O\coloneq \bigcup_{p+q=r} \Ad(G_0)o_{p,q}$.
	Then $O$ is a dense subset of $\lie{g}_1$ by Fact \ref{fact:JordanOrbitDecomposition}.

	Assume that $W$ is $\lie{g}_0$-stable.
	Then $W$ is an ideal of $\lie{g}_1$, i.e., $\jprod{\lie{g}_1, \lie{g}_1, W} \subset W$.
	This implies that the subalgebra generated by $W$ and $\theta(W)$ is a proper non-zero ideal of $\lie{g}$, which contradicts the simplicity of $\lie{g}$.
	Hence $W$ is not $\lie{g}_0$-stable.
	
	Since $W$ is not $\lie{g}_0$-stable, there exists $X \in W$ such that $[\lie{g}_0, X] \not\subset W$.
	Then we have $W + [\lie{g}_0, X] = \lie{g}_1$.
	Hence $\Ad(G_0)W$ contains an open neighborhood of $X$ in $\lie{g}_1$.
	Since $O$ is dense in $\lie{g}_1$, we have $\Ad(G_0)W \cap O \neq \emptyset$, which contradicts the assumption.
\end{proof}


\subsection{\texorpdfstring{$\lie{g} = \lieSO(p+1, q+1)$}{g=so(p+1,q+1)}} \label{subsect:JordanAlgebraSO(p,q)}

We shall consider the case of $\lie{g} = \lieSO(p+1, q+1)$ ($p, q \geq 0$, $p+q \geq 1$).
Retain the notation $\lie{g} = \lie{g}_{-1} \oplus \lie{g}_0 \oplus \lie{g}_1$, $\theta$, $h$ and $K_0$ from the previous subsection.

We have $\lie{g}_0 \simeq \lieSO(p, q) \oplus \RR h$ and $\lie{g}_1 \simeq \RR^{p, q}$ as a $\lie{g}_0$-module.
Here $\RR^{p, q}$ is the vector space $\RR^{p+q}$ equipped with the standard symmetric bilinear form $B$ of signature $(p, q)$.
Let $\langle \cdot, \cdot \rangle$ denote the standard inner product on $\RR^{p+q}$ and set
\begin{align*}
	I_{p,q} \coloneq \begin{pmatrix}
		I_p & 0 \\
		0 & -I_q
	\end{pmatrix}.
\end{align*}
Then we have $B(x,y) = \langle x, I_{p,q} y \rangle$.

The Jordan triple product is given by
\begin{align*}
	\jprod{x, y, z} &= \langle x, y \rangle z + \langle z, y \rangle x - B(x, z)I_{p,q}y \quad (x, y, z \in \RR^{p,q}).
\end{align*}
Let $\set{e_1, \ldots, e_{p+q}}$ be the standard basis of $\RR^{p,q}$ and set $u = e_1$.

\begin{proposition} \label{prop:TripotentSO(p,q)}
	Let $0\neq x \in \RR^{p,q}$.
	Then $x$ is a tripotent if and only if it satisfies exactly one of the following conditions:
	\begin{enumparen}
		\item $B(x, x) = 0$ and $\langle x, x \rangle = 1/2$,
		\item $\langle x, x \rangle = 1$ and $I_{p,q}x = x$,
		\item $\langle x, x \rangle = 1$ and $I_{p,q}x = -x$.
	\end{enumparen}
	Moreover, $x$ is unitary if and only if either (2) or (3) holds.
\end{proposition}

\begin{proof}
	The `if' part is verified by direct calculations.
	Assume that $x$ is a tripotent.
	Then we have $x = \jprod{x, x, x} = 2\langle x, x \rangle x - B(x, x)I_{p,q}x$.
	Hence $\set{x, I_{p,q}x}$ is linearly dependent or $B(x, x)=0$.
	The assertion follows from this.
\end{proof}

By Proposition \ref{prop:TripotentSO(p,q)}, any unitary tripotent is conjugate to one of $e_1$, $-e_1$, $e_{p+1}$ and $-e_{p+1}$ under the $K_0$-action.
Let $u \in \set{\pm e_1, \pm e_{p+1}}$, and set $\varepsilon = B(u, u)$.
The Cartan involution $\vartheta$ is given by
\begin{align*}
	\vartheta(x) = 2\langle x, u \rangle u - \varepsilon I_{p,q}x \quad (x \in \lie{g}_1).
\end{align*}
Hence, if $p, q\geq 1$, the set $\set{\frac{\varepsilon_1}{2}(e_1 + e_{p+1}), \frac{\varepsilon_2}{2}(e_1 - e_{p+1})}$ is a Jordan frame for some $\varepsilon_1, \varepsilon_2 \in \set{-1, 1}$.
If $p = 0$ or $q = 0$, then $\set{u}$ is a Jordan frame.

\begin{lemma} \label{lem:InvolutionSO(1,k)}
	One has
	\begin{align*}
		(\RR^{p, q})^{\vartheta} = \begin{cases}
			\spn{\RR}{e_1, e_{p+1}, \ldots, e_{p+q}} & (u=\pm e_1), \\
			\spn{\RR}{e_1, \ldots, e_p, e_{p+1}} & (u=\pm e_{p+1}),
		\end{cases} \\
		(\RR^{p,q})^{-\vartheta} = \begin{cases}
			\spn{\RR}{e_2, \ldots, e_p} & (u=\pm e_1), \\
			\spn{\RR}{e_{p+2}, \ldots, e_{p+q}} & (u=\pm e_{p+1}).
		\end{cases}
	\end{align*}
\end{lemma}

\subsection{Embedding of \texorpdfstring{$\lie{h} = \lieSO(1, k)$}{h=so(1,k)}}

In this subsection, we complete the proof of Theorem \ref{thm:MainTargetMinimal} in the case of type $A_1$ pairs.
Let $(\lie{g}, \lie{h})$ be a pair of type $A_1$ defined in Definition \ref{def:typeA}.
Then the embedding $\lie{h}_1 \rightarrow \lie{g}_1$ is a homomorphism of Jordan triple systems.
We remark that, in contrast to the case of type $BC_1$, $\lie{g}$ need not be simple.

Fix unitary tripotents $u' \in \lie{h}_1$ and $u \in \lie{g}_1$.
By Definition \ref{def:typeA}, $u'$ is unitary in $\lie{g}_1$.
Take a Jordan frame $\set{u_1, \ldots, u_r}$ in $(\lie{g}_1, u)^\vartheta$ and set $\lie{c}\coloneq \spn{\RR}{u_1, \ldots, u_r}$.
By Proposition \ref{prop:TripotentAndJordanFrame}, replacing $\lie{h}$ with a conjugate by an element of $K_0$, we may assume that $u' \in \lie{c}$.
Let $\vartheta'$ be the Cartan involution of $(\lie{g}_1, u')$ defined by $\vartheta'(x) = \jprod{u', x, u'}$.
Then the Jordan algebra $(\lie{h}_1, u')$ is a $\vartheta'$-stable subalgebra of $(\lie{g}_1, u')$.

Throughout this subsection, orthogonal complements are taken with respect to the form $-(\cdot,\theta(\cdot))$.

\begin{lemma} \label{lem:HperpInC}
	One has $\lie{c}\cap \lie{h} = \RR u'$ and $\lie{c} \cap \lie{h}^\perp = \lie{c} \cap (\RR u')^\perp$.
	In particular, one has $\lie{c} = (\lie{c}\cap \lie{h}) \oplus (\lie{c} \cap \lie{h}^\perp)$.
\end{lemma}

\begin{proof}
	By Lemma \ref{lem:InvolutionSO(1,k)}, we have $\lie{h}_1^{\vartheta'} = \RR u'$.
	By Proposition \ref{prop:TripotentAndJordanFrame}, there exist $\varepsilon_1, \ldots, \varepsilon_r \in \set{\pm 1}$ such that $u' = \sum_{i=1}^r \varepsilon_i u_i$.
	By Fact \ref{fact:JordanFrame}, we have $\lie{c} \subset \lie{g}_1^{\vartheta'}$ .
	Hence we obtain $\lie{c}\cap \lie{h} = \RR u'$.

	Since $\lie{g}_1^{\vartheta'}$ is orthogonal to $\lie{g}_1^{-\vartheta'}$, we have
	\begin{align*}
		\lie{c} \cap (\RR u')^\perp = \lie{c} \cap (\RR u')^\perp \cap (\lie{h}_1^{-\vartheta'})^\perp = \lie{c} \cap \lie{h}^\perp.
	\end{align*}
	We have shown the assertion.
\end{proof}

\begin{lemma} \label{lem:NonTrivialInvolution}
	If $k \geq 3$, one has $\vartheta' \neq \id_{\lie{g}_1}$.
\end{lemma}

\begin{proof}
	By Lemma \ref{lem:InvolutionSO(1,k)}, since $k \geq 3$, we have $\vartheta' \neq \id_{\lie{g}_1}$.
\end{proof}

First, we consider the case of simple $\lie{g}$.
Suppose that $\lie{g}_1^\vartheta$ is a simple Jordan algebra.
Define $K_0$, $G_0$, $o_{p,q}$ and $\lie{a'}$ as in Subsection \ref{subsect:JordanAlgebra}.
Then we have $u' = o_{a, r - a}$ ($0\leq a \leq r$).

\begin{lemma} \label{lem:NontrivialTripotent}
	If $\Delta(\lie{g}^{\theta\sigma}, \lie{a'})$ is of type $A_{r-1}$ and $k \geq 3$, then one has $a \neq 0, r$.
\end{lemma}

\begin{proof}
	Assume that $\Delta(\lie{g}^{\theta\sigma}, \lie{a'})$ is of type $A_{r-1}$.
	By Proposition \ref{prop:EuclideanJordanRoot}, we have $\vartheta = \id_{\lie{g}_1} \neq \vartheta'$.
	This implies that $u \neq \pm u'$ and hence $a \neq 0, r$.
\end{proof}

Let $(\lie{g}, \lie{h}, \orbit{O})$ be an NDD tuple.
By Lemma \ref{lem:ExistenceWeightVector}, there exists a non-zero element $X \in \lie{g}_1 \cap \lie{h}^\perp \cap \overline{\orbit{O}}$.

\begin{lemma} \label{lem:NonMinimalRank3}
	If $r \geq 3$, then $(\lie{g}, \lie{h}, \orbit{O})$ is not minimal.
\end{lemma}

\begin{proof}
	If $X \in \Ad(G_0)o_{1,0}\cup \Ad(G_0)o_{0,1}$, then $\set{[X, -\theta(X)], X, -\theta(X)}$ spans a subalgebra isomorphic to $\lieSL(2,\RR)$ by Lemma \ref{lem:RankOneTripotent}.
	Hence $(\lie{g}, \lie{h}, \orbit{O})$ is not minimal by Lemma \ref{lem:NonMinimalX}.
	
	Assume that $X \in \Ad(G_0)o_{p, q}$ with $p + q \geq 2$.
	We shall show $\lie{c}\cap \lie{h}^\perp \cap \overline{\orbit{O}} \neq \set{0}$, which is one of the assumptions in Lemma \ref{lem:ReductionRealStronglyOrthogonal}.
	The other assumption in Lemma \ref{lem:ReductionRealStronglyOrthogonal} is verified in Lemma \ref{lem:HperpInC}.
	Note that $\lie{g}$ is not isomorphic to $\lieSL(2,\RR)$.

	Assume that $\Delta(\lie{g}^{\theta\sigma}, \lie{a'})$ is of type $A_{r-1}$ and $u' = \pm u$.
	By Fact \ref{fact:RootSystemJordan}, $K_0$ fixes $u'$.
	Take $k \in K_0$ such that $\Ad(k)X \in \lie{c}$ by Proposition \ref{prop:K0Orbit}.
	Then we have $(\Ad(k)X, -\theta(u')) = (X, -\theta(u')) = 0$.
	By Lemma \ref{lem:HperpInC}, we obtain $\Ad(k)X \in \lie{c} \cap \lie{h}^\perp$, and hence $\lie{c}\cap \lie{h}^\perp \cap \overline{\orbit{O}} \neq \set{0}$.

	Assume that $\Delta(\lie{g}^{\theta\sigma}, \lie{a'})$ is not of type $A_{r-1}$ or $u' \neq \pm u$.
	By Fact \ref{fact:WeylGroupJordan}, if necessary, replacing $u'$ by $\varepsilon wu'$ for some $w \in W'$, $\varepsilon \in \set{-1, 1}$, we may assume that $u'=o_{a,r-a}$ ($a\geq r-a > 0$), and then $a \geq 2$.
	Since $X \in \Ad(G_0)o_{p, q}$ with $p + q \geq 2$, $\overline{\orbit{O}}$ contains one of the following tripotents:
	\begin{align*}
		u_a+u_{a+1}, \quad -u_a-u_{a+1}, \quad u_1-u_2
	\end{align*}
	by Fact \ref{fact:WeylGroupJordan}.
	The elements belong to $\lie{c} \cap (\RR u')^\perp$ and hence to $\lie{h}^\perp$ by Lemma \ref{lem:HperpInC}.
	Therefore we obtain $\lie{c}\cap \lie{h}^\perp \cap \overline{\orbit{O}} \neq \set{0}$.
\end{proof}

\begin{lemma}
	If $\Delta(\lie{g}^{\theta\sigma}, \lie{a'})$ is of type $C_{r}$ and $r\geq 2$, then $(\lie{g}, \lie{h}, \orbit{O})$ is not minimal.
\end{lemma}

\begin{proof}
	Under the assumptions, in the proof of Lemma \ref{lem:NonMinimalRank3}, we may choose the element in $\overline{\orbit{O}}$ to be $u_1 + u_2$ using the Weyl group action.
	Then we have $u_1 + u_2 \in \lie{c} \cap (\RR u')^\perp = \lie{c} \cap \lie{h}^\perp$ by Lemma \ref{lem:HperpInC}.
	This shows the lemma.
\end{proof}

The remaining cases are those in which $\Delta(\lie{g}^{\theta\sigma}, \lie{a'})$ is of type $A_{r-1}$ or $D_r$ ($r \leq 2$).
By the classification (Table \ref{table:JordanAlgebra}), the only possibility is the case of $\lie{g} = \lieSO(p+1, q+1)$ ($p, q \geq 1$).
Note that, if $p = 0$ or $q = 0$, then $(\lie{g}, \lie{h}, \orbit{O})$ is not minimal by Corollary \ref{cor:SufficientSameRank}.
Then we have $r = 2$.

As in Subsection \ref{subsect:JordanAlgebraSO(p,q)}, identify $\lie{g}_1$ with $\RR^{p,q}$.
Retain the notation $I_{p,q}$, $B$ and $\langle\cdot, \cdot \rangle$ from Subsection \ref{subsect:JordanAlgebraSO(p,q)}.
By Proposition \ref{prop:TripotentSO(p,q)}, swapping $p$ and $q$ if necessary, we may assume that $u = u' = e_1$, and $u_1 = (e_1 + e_{p+1})/2, u_2 = (e_1 - e_{p+1})/2$.
Then we have $\vartheta' = \vartheta$.
Here, we do not assume that $\lie{g}_1^{\vartheta}$ is simple.

\begin{lemma}\label{lem:Rank2Orthogonal}
	One has $\lie{g}_1^{-\vartheta} = \lie{h}_1^{-\vartheta}$.
\end{lemma}

\begin{proof}
	Assume that $\lie{g}_1^{-\vartheta} \neq \lie{h}_1^{-\vartheta}$.
	Recall that the bilinear form $B$ is definite on $\lie{g}_1^{-\vartheta}$ by Lemma \ref{lem:InvolutionSO(1,k)}.
	By the action of $K_0$, we may assume that $\lie{h}_1 = \spn{\RR}{e_1, \ldots, e_{k-1}}$.
	Then we have
	\begin{align*}
		\lie{g}_1 \cap \lie{h}^\perp = \spn{\RR}{e_k, \ldots, e_{p+q}}.
	\end{align*}
	This implies that the bilinear form $B$ is indefinite on $\lie{g}_1 \cap \lie{h}^\perp$.
	Let $\lie{h'}$ be the subalgebra of $[\lie{g}_0, \lie{g}_0] \simeq \lieSO(p, q)$ isomorphic to $\lieSO(p - (k-1), q)$.
	This subalgebra is the Lie algebra of the isometry group of $\lie{g}_1\cap \lie{h}^\perp$.
	Then $\lie{h'}$ is $\theta$-stable and acts on $\lie{h}_1$ trivially.
	Since $\lie{h}$ is generated by $\lie{h}_1$ and $\theta(\lie{h}_1)$, we have $\lie{h'}\subset \lieCent_{\lie{g}}(\lie{h})$.
	This contradicts the assumption that $\lieCent_{\lie{g}}(\lie{h})$ is compact.
\end{proof}

\begin{lemma}
	Under the above assumptions, $(\lie{g}, \lie{h}, \orbit{O})$ is not minimal.
\end{lemma}

\begin{proof}
	By Lemma \ref{lem:Rank2Orthogonal}, we have $p+1 = k$ and $\lieNorm_{\lie{g}}(\lie{h}) \simeq \lieSO(k, 1) \oplus \lieSO(0, q)$.
	This implies that $(\lie{g}, \lieNorm_{\lie{g}}(\lie{h}))$ is a symmetric pair.
	By Lemma \ref{lem:SufficientSymmetric}, $(\lie{g}, \lie{h}, \orbit{O})$ is not minimal.
\end{proof}

\begin{theorem} \label{thm:NonMinimalTypeASimple}
	Let $(\lie{g}, \lie{h}, \orbit{O})$ be an NDD tuple such that $(\lie{g}, \lie{h})$ is of type $A_1$.
	If $\lie{g}$ is simple, then $(\lie{g}, \lie{h}, \orbit{O})$ is not minimal.
\end{theorem}

Next, we shall consider the case in which $\lie{g}$ is not simple.
Let $\lie{g} = \lie{g}^1\oplus \lie{g}^2 \oplus \cdots \oplus \lie{g}^m$ ($m \geq 2$) be the decomposition of $\lie{g}$ into simple ideals.
Then $\lie{g}_1 = \lie{g}_1^1 \oplus \lie{g}_1^2 \oplus \cdots \oplus \lie{g}_1^m$ is the decomposition of $\lie{g}_1$ into simple Jordan triple systems.
Write $p_i\colon \lie{g} \rightarrow \lie{g}^i$ for the projection onto the $i$-th component.
By Definition \ref{def:typeA}, we have $\lie{g}_1^i \neq 0$ and $p_i(\lie{h}) \neq 0$ for any $i$.

\begin{lemma} \label{lem:SL2DirectSum}
	Assume that $\lie{h} \simeq \lieSL(2,\RR)$ and every simple factor of $\lie{g}$ is isomorphic to $\lieSL(2,\RR)$.
	Then $(\lie{g}, \lie{h}, \orbit{O})$ is not minimal for any $\orbit{O}$.
\end{lemma}

\begin{proof}
	We may assume that $\lie{g} = \lieSL(2,\RR)^{\oplus m}$ and $\lie{h}$ is the diagonal subalgebra of $\lie{g}$.
	As a $G$-invariant bilinear form on $\lie{g}$, we take the Killing form.
	Let $\set{h_0, x_0, y_0}$ be an $\lieSL_2$-triple in $\lieSL(2,\RR)$ such that $2h = (h_0, h_0, \ldots, h_0)$.
	By Lemma \ref{lem:ExistenceWeightVector}, $\lie{g}_1 \cap \lie{h}^\perp \cap \overline{\orbit{O}}$ contains
	a nonzero element $X = (a_1 x_0, a_2 x_0, \ldots, a_m x_0)$ for some $a_i \in \RR$.
	Since $X \in \lie{h}^\perp$, we have $\sum_{i=1}^m a_i = 0$.
	In particular, there exist $i, j$ such that $a_i > 0$ and $a_j < 0$.
	We may assume that $a_1 > 0$ and $a_2 < 0$.

	By acting with $G_0$ and passing to a suitable limit, we have
	\begin{align*}
		X' \coloneq (x_0, -x_0, 0, \ldots, 0) \in \lie{g}_1\cap \overline{\orbit{O}}\cap \lie{h}^\perp.
	\end{align*}
	Since $\set{[X', -\theta(X')], X', -\theta(X')}$ is an $\lieSL_2$-triple, $(\lie{g}, \lie{h}, \orbit{O})$ is not minimal.
\end{proof}

\begin{theorem} \label{thm:NonMinimalTypeASemisimple}
	Let $(\lie{g}, \lie{h}, \orbit{O})$ be an NDD tuple such that $(\lie{g}, \lie{h})$ is of type $A_1$.
	If $\lie{g}$ is not simple, then $(\lie{g}, \lie{h}, \orbit{O})$ is not minimal.
\end{theorem}

\begin{proof}
	By Lemma \ref{lem:ExistenceWeightVector}, there exists a non-zero element $X \in \lie{g}_1 \cap \lie{h}^\perp \cap \overline{\orbit{O}}$.
	If $p_i(X) = 0$ for some $i$, then $(\lie{g}, \lie{h}, \orbit{O})$ is not minimal by Proposition \ref{prop:BasicMinimal} (6).
	Hence we may assume that $p_i(X) \neq 0$ for any $i$.

	First, assume that $k \geq 3$, i.e., $\lie{h}$ is not isomorphic to $\lieSL(2,\RR)$.
	Note that $\lie{g}^i$ is not isomorphic to $\lieSL(2,\RR)$ for any $i$ since $p_i(\lie{h}) \neq 0$.
	By Lemma \ref{lem:NontrivialTripotent}, the assumption of Lemma \ref{lem:JordanAnyInnerProduct} is satisfied for each simple factor.
	Hence we can take an element $X' \in \lie{c} \cap \overline{\orbit{O}}$ such that $(X', -\theta(u')) = 0$.
	By Lemma \ref{lem:HperpInC}, this implies that $X' \in \lie{h}^\perp$.
	In particular, we have  $\lie{c} \cap \lie{h}^\perp \cap \overline{\orbit{O}} \neq \set{0}$.
	By Lemma \ref{lem:HperpInC}, again, we have $\lie{c} = (\lie{c}\cap \lie{h}) \oplus (\lie{c} \cap \lie{h}^\perp)$.
	Therefore, $\lie{c}$ satisfies the assumptions of Lemma \ref{lem:ReductionRealStronglyOrthogonal}, and hence $(\lie{g}, \lie{h}, \orbit{O})$ is not minimal.

	Next, assume that $k = 2$, i.e., $\lie{h} \simeq \lieSL(2,\RR)$.
	If every simple factor of $\lie{g}$ is isomorphic to $\lieSL(2,\RR)$, then $(\lie{g}, \lie{h}, \orbit{O})$ is not minimal by Lemma \ref{lem:SL2DirectSum}.
	If the involution $\vartheta'(\cdot) = \jprod{u', \cdot, u'}$ is not the identity, $(\lie{g}, \lie{h}, \orbit{O})$ is not minimal by the same argument as in the first case.
	
	Hence we may assume that $\lie{g}$ has a simple factor not isomorphic to $\lieSL(2,\RR)$, and $\jprod{u', \cdot, u'}$ is the identity, which implies that every simple factor of $\lie{g}_1$ is Euclidean by Proposition \ref{prop:EuclideanJordanRoot}.
	By Fact \ref{fact:RootSystemJordan}, $u'$ is a $K_0$-fixed vector.
	Take $k \in K_0$ such that $\Ad(k)X \in \lie{c}$ by Proposition \ref{prop:K0Orbit}.
	Then we have $(\Ad(k)X, -\theta(u')) = (X, -\theta(u')) = 0$.
	This implies that $\Ad(k)X \in \lie{c} \cap \lie{h}^\perp$, and hence $\lie{c}\cap \lie{h}^\perp \cap \overline{\orbit{O}} \neq \set{0}$.
	Therefore, $\lie{c}$ satisfies the assumptions of Lemma \ref{lem:ReductionRealStronglyOrthogonal}, and hence $(\lie{g}, \lie{h}, \orbit{O})$ is not minimal.
\end{proof}